% !PRINTED 26
\weberpage{26}

%<<CONT
occurs, and that the indices in this part of the series are all
equal to $\nu$.

The assertion is true for $r = 2$, and we prove it in general
by complete induction, under the assumption that it has already
been proved true for $r$.

If $Q_r$ is not equal to $R$, we choose $Q_1^{(r)}$ so that $Q_{r+1}$ is the
intersection of $Q_1, Q_1' \cdots Q_1^{(r-1)}, Q_1^{(r)}$, and moreover denote
the intersection of $Q_1, Q_1' \cdots Q_1^{(r-2)}, Q_1^{(r)}$ by $Q_r'$. Then
$Q_r'$ is distinct from $Q_r$, and $Q_{r+1}$ is the intersection of
$Q_r$ and $Q_r'$. Now by the assumption (4) we have two
corresponding pieces of the composition series with the constant
index series $\nu$

\weq{5}{
\begin{array}{l}
Q,\ Q_1 \cdots Q_{r-1}\ Q_r,\\
Q,\ Q_1 \cdots Q_{r-1}\ Q_r';
\end{array}
}

for the second of these series is formed exactly like the first
according to the prescription (3), only that $Q_1^{(r)}$ has taken the
place of $Q_1^{(r-1)}$. By the theorem of \S.\,6, 2., therefore, $Q_{r+1}$ is
the greatest normal divisor both of $Q_r$ and of $Q_r'$, and indeed
of index $\nu$, and it follows therefore, as was to be proved, that
the composition series can be continued thus:

\weq{6}{Q,\ Q_1 \cdots Q_{r-1},\ Q_r,\ Q_{r+1},}

and that the newly added index is likewise $\nu$.

This we can now continue until we have reached the group $R$,
and we obtain a piece of the composition series

\weq{7}{
\begin{array}{ccc}
Q,\ Q_1,\ Q_2 & \cdots & R,\\
\nu\quad\ \nu & & \nu
\end{array}
}

whose last term $R$ is again a normal divisor of $P$ and in
which all the indices agree.

If we continue the composition series onward from $R$, the
index may change; but from there on it can again be kept
equal, until one once more arrives at a normal divisor of $P$.

With this the following theorem is proved:

III.\ \gesperrt{One can so arrange the composition series of $P$ with
its index series}

\weq{8}{
\begin{array}{c}
P,\ P_1,\ P_2,\ P_3 \cdots\\
j_1,\ j_2,\ j_3 \cdots
\end{array}
}

\gesperrt{that, whenever a change occurs in the index
series, that is, whenever $j_{r+1}$ is different from $j_r$,
$P_r$ is a normal divisor of $P$.}
